% GENERATED FILE. Edit canonical lessons, then run npm run generate:books.
% canonical-source-hash: fnv1a64:e9fe7e068bed0314
% canonical-lessons: ZH-C474-gebo, ZH-C474-liqi, ZH-W474-li-strength, ZH-C474-tili, ZH-C474-yongli, ZH-C474-shouzhi

\chapter{Strength}
\label{ch:zh-strength}
\hlchaptermodality{474}

\begin{chapteropening}
\textbf{By the end of this chapter:} \emph{I can talk about my arms and fingers, how strong someone is, their stamina, and pressing hard.}

Complete the last lesson of chapter 474: I can talk about my arms and fingers, how strong someone is, their stamina, and pressing hard.
\end{chapteropening}

\section[gēbo]{gēbo — arm}

\label{lesson:ZH-C474-gebo}



\begin{quote}
Before the new one: say the Mandarin for to have a tooth pulled, then the Mandarin for to have a tooth filled.
\end{quote}

\subsection*{You'll want to know: gēbo}
\emph{gēbo} — \textquotedblleft{}arm\textquotedblright{}.

\begin{quote}
\emph{wǒ gōngzuò le shí ge xiǎoshí gēbo hěn téng} — I worked for ten hours, and my arms really ache.
\end{quote}

Its characters come later; for now it is a word for the ear.

\begin{cousinweb}
\emph{gēbo} is built from \emph{gē} and \emph{bo}, both inherited from Middle Chinese, \emph{kak} and \emph{phak} (Guangyun, in Baxter's transcription). The dictionary glosses the first as the armpit and the second as dried meat; how the pair came to mean the arm is not settled, and the first tone of \emph{gē} is irregular.
\end{cousinweb}

\begin{grammarlens}[title={What you've built}]
One more word for saying how it was.
\end{grammarlens}

\subsection*{Your turn}
\begin{itemize}
\raggedright
  \item \emph{Say it:} \emph{gēbo}
  \item \emph{Say it:} \emph{gēbo}, once more
  \item \emph{Say it:} \emph{wǒ gōngzuò le shí ge xiǎoshí gēbo hěn téng} — I worked for ten hours, and my arms really ache.
  \item \emph{Recall:} say the Mandarin for to have a tooth pulled, then the Mandarin for to have a tooth filled, then say \emph{gēbo} again
\end{itemize}

\begin{tcolorbox}[breakable,colback=teal!4,colframe=teal!35!black,title={Before you move on}]
What does gēbo mean? (\textquotedblleft{}Arm\textquotedblright{}.) Say it once more.
\end{tcolorbox}
\section[\texorpdfstring{\zh{力气} (lìqi)}{力气 (lìqi)}]{\zh{力气} (lìqi) — physical strength}

\label{lesson:ZH-C474-liqi}



\begin{quote}
Before the new one: say the Mandarin for to have a tooth filled, then the Mandarin for the arm.
\end{quote}

\subsection*{You'll want to know: \zh{力气}}
\textbf{\zh{力气}} — \emph{lìqi} — \textquotedblleft{}physical strength\textquotedblright{}.

\begin{quote}
\textbf{\zh{他的力气比我大}} \emph{tā de lìqi bǐ wǒ dà} — He's stronger than I am.
\end{quote}

The new character, in pieces: \zh{力} — two strokes: a stroke that runs right, turns down and hooks to the left, and a long stroke falling to the left through it. The next lesson puts it in your hand.

\begin{cousinweb}
\textbf{\zh{力气}} (\emph{lìqi}) is built from \zh{力} 'strength' and \zh{气} 'breath', both inherited from Middle Chinese, \emph{lik} and \emph{khj+jH} (Guangyun, in Baxter's transcription), which glosses the first 'sinew' and the second 'breath'; the second is said light.
\end{cousinweb}

\begin{grammarlens}[title={What you've built}]
One more word for saying how it was.
\end{grammarlens}

\subsection*{Your turn}
\begin{itemize}
\raggedright
  \item \emph{Say it:} \emph{lìqi}
  \item \emph{Say it:} \emph{lìqi}, once more
  \item \emph{Say it:} \emph{tā de lìqi bǐ wǒ dà} — He's stronger than I am.
  \item \emph{Recall:} say the Mandarin for to have a tooth filled, then the Mandarin for the arm, then say \emph{lìqi} again
\end{itemize}

\begin{tcolorbox}[breakable,colback=teal!4,colframe=teal!35!black,title={Before you move on}]
What does \zh{力气} mean? (\textquotedblleft{}Physical strength\textquotedblright{}.) Say it once more.
\end{tcolorbox}
\section[\texorpdfstring{\zh{力} (lì)}{力 (lì)}]{\zh{力} — a turn with a hook, then a falling stroke}

\label{lesson:ZH-W474-li-strength}



\begin{quote}
The word first, then the character inside it.

\begin{itemize}
\raggedright
  \item \emph{Recall:} say \emph{lìqi} — \textbf{R1}, one lesson back, before it decays
\end{itemize}

\begin{quote}
\zh{男} \zh{气}
\end{quote}

\emph{Read it:} each one aloud, with nothing said first — \emph{nán}, \emph{qì}
\end{quote}

\subsection*{Script — a turn with a hook, then a falling stroke}
\hlblockfigure{\detokenize{figures/ZH-W474-li-strength-filmstrip.pdf}}{How \zh{力} is written, stroke by stroke}

\begin{quote}
\zh{力}
\end{quote}

\textbf{One piece of two strokes.}

\textbf{The first stroke.} It runs right, turns down, and at the foot hooks back to the left.

\textbf{The second stroke.} A long stroke falling to the left, from above the first stroke, down through it to the bottom left.

1 + 1 = two.

\subsection*{Your turn}
\begin{itemize}
\raggedright
  \item \emph{Write it:} \zh{力}, all two, with the model in front of you
  \item \emph{Cover:} write it again from memory, and count the pieces, not the strokes
\end{itemize}

\begin{tcolorbox}[breakable,colback=teal!4,colframe=teal!35!black,title={Before you move on}]
How many strokes does \zh{力} have? (\textbf{Two}.) Where have you written it before? (\textbf{At the foot of \zh{男}}.)

Source: \href{https://raw.githubusercontent.com/chanind/hanzi-writer-data/68d10a4b21150cae5e1ebbd223eed289cf32d90c/data/\%E5\%8A\%9B.json}{Hanzi Writer Data \zh{力}, pinned stroke paths and medians}.
\end{tcolorbox}
\section[\texorpdfstring{\zh{体力} (tǐlì)}{体力 (tǐlì)}]{\zh{体力} (tǐlì) — stamina; physical fitness}

\label{lesson:ZH-C474-tili}



\begin{quote}
Before the new one: say the Mandarin for the arm, then the Mandarin for physical strength.
\end{quote}

\subsection*{You'll want to know: \zh{体力}}
\textbf{\zh{体力}} — \emph{tǐlì} — \textquotedblleft{}stamina; physical fitness\textquotedblright{}.

\begin{quote}
\textbf{\zh{他病了一个月脸色和体力都不好}} \emph{tā bìng le yí ge yuè liǎnsè hé tǐlì dōu bù hǎo} — He was ill for a month, and his colour and stamina are both poor.
\end{quote}

\begin{cousinweb}
\textbf{\zh{体力}} (\emph{tǐlì}) is built from \zh{体} 'body' and \zh{力} 'strength', both inherited from Middle Chinese, \emph{thejX} and \emph{lik} (Guangyun, in Baxter's transcription), which glosses the first 'the body' and the second 'sinew'.
\end{cousinweb}

\begin{grammarlens}[title={What you've built}]
One more word for saying how it was.
\end{grammarlens}

\subsection*{Your turn}
\begin{itemize}
\raggedright
  \item \emph{Say it:} \emph{tǐlì}
  \item \emph{Say it:} \emph{tǐlì}, once more
  \item \emph{Say it:} \emph{tā bìng le yí ge yuè liǎnsè hé tǐlì dōu bù hǎo} — He was ill for a month, and his colour and stamina are both poor.
  \item \emph{Recall:} say the Mandarin for the arm, then the Mandarin for physical strength, then say \emph{tǐlì} again
\end{itemize}

\begin{tcolorbox}[breakable,colback=teal!4,colframe=teal!35!black,title={Before you move on}]
What does \zh{体力} mean? (\textquotedblleft{}Stamina; physical fitness\textquotedblright{}.) Say it once more.
\end{tcolorbox}
\section[\texorpdfstring{\zh{用力} (yònglì)}{用力 (yònglì)}]{\zh{用力} (yònglì) — to push hard; to put one's strength into it}

\label{lesson:ZH-C474-yongli}



\begin{quote}
Before the new one: say the Mandarin for physical strength, then the Mandarin for stamina.
\end{quote}

\subsection*{You'll want to know: \zh{用力}}
\textbf{\zh{用力}} — \emph{yònglì} — \textquotedblleft{}to push hard; to put one's strength into it\textquotedblright{}.

\begin{quote}
\textbf{\zh{牙医说吃东西的时候不要太用力}} \emph{yáyī shuō chī dōngxi de shíhou bú yào tài yònglì} — The dentist says not to bite down too hard when I eat.
\end{quote}

\begin{cousinweb}
\textbf{\zh{用力}} (\emph{yònglì}) is built from \zh{用} 'to use' and \zh{力} 'strength', both inherited from Middle Chinese, \emph{yowngH} and \emph{lik} (Guangyun, in Baxter's transcription).
\end{cousinweb}

\begin{grammarlens}[title={What you've built}]
One more word for saying how it was.
\end{grammarlens}

\subsection*{Your turn}
\begin{itemize}
\raggedright
  \item \emph{Say it:} \emph{yònglì}
  \item \emph{Say it:} \emph{yònglì}, once more
  \item \emph{Say it:} \emph{yáyī shuō chī dōngxi de shíhou bú yào tài yònglì} — The dentist says not to bite down too hard when I eat.
  \item \emph{Recall:} say the Mandarin for physical strength, then the Mandarin for stamina, then say \emph{yònglì} again
\end{itemize}

\begin{tcolorbox}[breakable,colback=teal!4,colframe=teal!35!black,title={Before you move on}]
What does \zh{用力} mean? (\textquotedblleft{}To push hard; to put one's strength into it\textquotedblright{}.) Say it once more.
\end{tcolorbox}
\section[shǒuzhǐ]{shǒuzhǐ — finger}

\label{lesson:ZH-C474-shouzhi}



\begin{quote}
Before the new one: say the Mandarin for stamina, then the Mandarin for to put one's strength into it.
\end{quote}

\subsection*{You'll want to know: shǒuzhǐ}
\emph{shǒuzhǐ} — \textquotedblleft{}finger\textquotedblright{}.

\begin{quote}
\emph{wǒ gǎnjué shǒuzhǐ hěn lěng} — My fingers feel really cold.
\end{quote}

Its characters come later; for now it is a word for the ear.

Two third tones meet in this word: shǒu is said with a rising tone before zhǐ, and the spelling keeps both marks.

\begin{cousinweb}
\emph{shǒuzhǐ} is built from \zh{手} 'hand' and \emph{zhǐ} 'finger', both inherited from Middle Chinese, \emph{syuwX} and \emph{tsyijX} (Guangyun, in Baxter's transcription), which glosses the second with this very word.
\end{cousinweb}

\begin{grammarlens}[title={What you've built}]
One more word for saying how it was.
\end{grammarlens}

\subsection*{Your turn}
\begin{itemize}
\raggedright
  \item \emph{Say it:} \emph{shǒuzhǐ}
  \item \emph{Say it:} \emph{shǒuzhǐ}, once more
  \item \emph{Say it:} \emph{wǒ gǎnjué shǒuzhǐ hěn lěng} — My fingers feel really cold.
  \item \emph{Recall:} say the Mandarin for stamina, then the Mandarin for to put one's strength into it, then say \emph{shǒuzhǐ} again
\end{itemize}

\begin{tcolorbox}[breakable,colback=teal!4,colframe=teal!35!black,title={Before you move on}]
What does shǒuzhǐ mean? (\textquotedblleft{}Finger\textquotedblright{}.) Say it once more.
\end{tcolorbox}
